\section{Centrales comerciales: financiamiento, riesgos y elección de ruta}

Hasta ahora se habló de los dispositivos; esta sección trata del dinero y de la ruta. La ruta que plantea Yang Zhao (杨钊) se divide en tres dispositivos clave: Honghuang 70 (洪荒 70) verifica la viabilidad de ingeniería de un dispositivo completo; Honghuang 170 alcanza parámetros altos y una Q mayor que 10 con un costo relativamente bajo; Honghuang 380 funciona como central de demostración completa, con el objetivo de convertirse en un producto que pueda venderse a los propietarios de centrales eléctricas. El 170 implica una inversión del orden de tres mil millones de yuanes, mientras que el 380 corresponde a un costo de central de demostración del orden de diez mil a veinte mil millones de yuanes.

Lo más importante aquí es la transformación del riesgo. Nengliang Qidian (能量奇点) elige un diseño físico relativamente conservador y procura no asumir nuevos riesgos científicos; lo que en realidad debe resolver son los riesgos de ingeniería de los parámetros altos y los riesgos de costo. Frente a otras rutas como la de Helion, la ventaja del tokamak superconductor de alta temperatura es que la ruta del tokamak ya cuenta con una gran acumulación de experimentos de parámetros altos; en cambio, algunas otras configuraciones de confinamiento magnético quizá todavía necesiten extrapolar sus parámetros experimentales varios órdenes de magnitud, lo cual, a juicio de Yang Zhao, constituye una ruta con mayor riesgo científico.

\begin{warningbox}{Hay que separar el riesgo científico del riesgo de ingeniería}
El riesgo científico es «si la respuesta existe y si los procesos físicos funcionan según la extrapolación»; el riesgo de ingeniería es «la respuesta existe, pero ¿se puede construir un sistema que cumpla con los parámetros, el costo y la confiabilidad?». Una empresa emergente es más adecuada para reducir el riesgo de ingeniería y el riesgo de costo; los problemas de alto riesgo científico son más adecuados para la exploración a largo plazo de las instituciones de investigación.
\end{warningbox}

\begin{knowledgebox}{Tabla de financiamiento e hitos}
\begin{tabularx}{\textwidth}{p{0.18\textwidth}p{0.32\textwidth}X}
\toprule
Etapa & Objetivo & Implicación financiera/comercial \\
\midrule
Honghuang 70 & Viabilidad de ingeniería de un dispositivo superconductor de alta temperatura completo & Demostrar que la ruta del nuevo material puede construirse y operar. \\
Honghuang 170 & Alcanzar Q mayor que 10 a bajo costo & Requiere un financiamiento de gran monto; demuestra que el dispositivo central para la comercialización es viable. \\
Honghuang 380 & Central de demostración; concepto de producto de unos 500MW & Venta a los propietarios de centrales eléctricas; verifica la rentabilidad de la central. \\
Replicación a escala & Múltiples centrales comerciales & El costo por kilowatt-hora sigue bajando y se entra a competir en el mercado energético. \\
\bottomrule
\end{tabularx}
\end{knowledgebox}

\begin{figure}[H]
\centering
\includegraphics[width=0.92\textwidth]{holy-grail-tradeoff.png}
\caption{El precio del santo grial de la energía: detrás de la visión de una energía ilimitada hay una complejidad de ingeniería extrema. Diagrama conceptual propio, elaborado a partir de la conversación de 01:30:00--02:00:00.}
\end{figure}

\begin{knowledgebox}{Lectura de la figura: el grial es difícil porque el costo aparece a la vez en cinco direcciones}
La alta temperatura, el confinamiento magnético, la vida útil de los materiales, el consumo de capital y la incertidumbre de los plazos conforman, en conjunto, la dificultad de comercializar la fusión. No basta con ser optimista en un solo punto: hay que poner la física, la ingeniería, el capital y la regulación en un mismo libro de cuentas.
\end{knowledgebox}

\subsection{El costo por kilowatt-hora es el criterio de referencia}

Esta sección reduce todas las rutas técnicas anteriores a un juicio comercial. El valor de Q, el campo magnético, el tamaño del dispositivo y los parámetros de los subsistemas son importantes, pero en última instancia deben servir a un indicador sencillo: cuánto cuesta generar un kilowatt-hora. Solo cuando el costo por kilowatt-hora se acerca al de la generación termoeléctrica, la fusión empieza a entrar al sistema energético; si llegara a ser un orden de magnitud menor, podría cambiar la forma en que se usa la energía.

En la entrevista, Yang Zhao insiste en que el criterio de referencia para comercializar la fusión es el costo por kilowatt-hora. Mientras ese costo se acerque al de la generación termoeléctrica, la fusión tiene posibilidades de entrar al sistema energético; si pudiera ser un orden de magnitud menor que el de la termoeléctrica, tanto el panorama del suministro energético como la forma en que la civilización usa la energía cambiarían más a fondo. La primera central de demostración tendrá forzosamente un costo elevado, pero si ya se acerca al de la termoeléctrica, el mercado podrá ver el margen de reducción que traerá la replicación a escala.

\begin{importantbox}{Comercializar no es el primer kilowatt-hora, sino el primer kilowatt-hora de bajo costo}
Generar el primer kilowatt-hora es un hito técnico; generar electricidad de forma estable con costos aceptables de construcción, regulación y operación y mantenimiento es el hito comercial. Entre la demostración y su conversión en una fuente de energía principal, la fusión nuclear todavía debe superar la operación prolongada, la extracción de calor para generar electricidad, el mantenimiento, la regulación y la replicación de la cadena de suministro.
\end{importantbox}

\begin{knowledgebox}{Implicaciones comerciales de los tres dispositivos}
\begin{tabularx}{\textwidth}{p{0.18\textwidth}p{0.34\textwidth}X}
\toprule
Dispositivo & Objetivo principal & Interpretación comercial \\
\midrule
Honghuang 70 & Verificar que un tokamak superconductor de alta temperatura completo pueda construirse y operar & Demuestra que «el barco de acero puede botarse al agua»; no demuestra directamente la rentabilidad de la central. \\
Honghuang 170 & Dispositivo experimental que alcance Q mayor que 10 con un costo relativamente bajo & Demuestra que la ruta superconductora de alta temperatura puede convertir un dispositivo científico en uno comercialmente viable. \\
Honghuang 380 & Concepto de central de demostración de unos 500MW, para entregarse a propietarios & Pone a prueba la operación prolongada, la extracción de calor para generar electricidad y la conversión de la central en producto. \\
\bottomrule
\end{tabularx}
\end{knowledgebox}

\subsection{Resumen de la sección}

El verdadero producto de una empresa emergente de fusión no son los artículos ni las fotografías de dispositivos, sino una central con un costo aceptable. La esencia de la ruta Honghuang consiste en adelantar al máximo el riesgo científico hasta convertirlo en física ya verificada, y concentrar los recursos de la empresa emergente en la entrega de ingeniería, la verificación de los subsistemas y la reducción del costo por kilowatt-hora.
